% !TEX program = xelatex
\documentclass[11pt]{letter}
\usepackage{fontspec}
\setmainfont{Times New Roman}
\usepackage{polyglossia}
\setdefaultlanguage{russian}
\usepackage[margin=1in]{geometry}
\usepackage{hyperref}
\setlength{\parindent}{0pt}
\setlength{\parskip}{0.6\baselineskip}
\frenchspacing

\signature{Никита Поломошнов}
\address{}
\date{18 апреля 2026 г.}

\begin{document}
\begin{letter}{[Имя Отчество]}
\opening{Здравствуйте, [Имя Отчество].}

Я хотел бы обсудить с Вами возможность научного руководства или консультации по проекту TGNN-Solv и, если это возможно, доступ к вычислительным ресурсам для завершения ключевой части экспериментов.

Проект посвящён предсказанию растворимости твёрдых органических веществ в растворителях. Я построил физически обоснованную графовую модель, которая не предсказывает растворимость напрямую, а проходит через уравнение твёрдо-жидкого равновесия: отдельно оценивает вклад плавления чистого вещества и вклад неидеальности раствора. Для честного сравнения реализована контрольная модель с тем же графовым представлением, но без физического слоя.

Сейчас уже готовы данные, код обучения, физический решатель, внешние сравнения, температурные диагностические разбиения и расширенная вспомогательная супервизия по коэффициентам активности. Основной текущий результат такой: на строгом разбиении по новым молекулярным каркасам прямая модель пока лучше физической модели по MAE $1.652$ против $1.741$. При этом температурные проверки показывают, что физический сигнал в данных очень силён: для одной и той же пары вещество--растворитель аппроксимация Вант-Гоффа даёт MAE $0.368$ на переносе с низких температур на высокие.

Поэтому следующий вопрос уже вычислительный: нужно проверить, сможет ли физическая модель при полном обучении использовать этот температурный сигнал лучше прямой модели. На локальной машине я могу проводить только короткие диагностические запуски; они полезны для отладки, но недостаточны для окончательного вывода.

Минимально полезный ресурс для следующего этапа --- одна CUDA-GPU с $24$ ГБ видеопамяти примерно на $300$ GPU-часов. Для результата, пригодного для публикации, лучше иметь $800$--$1000$ GPU-часов, чтобы провести несколько повторов и абляционные проверки.

Я приложил краткий отчёт и более полную версию с таблицами, математическими выкладками и диагностикой. Буду благодарен, если Вы сможете посмотреть, насколько проект выглядит научно состоятельным и есть ли возможность обсудить доступ к вычислительным ресурсам.

\closing{С уважением,}
\end{letter}
\end{document}
